\documentclass[12pt,a4paper]{article}
\usepackage[utf8]{inputenc}
\usepackage{amsmath}
\usepackage{graphicx}
\usepackage{hyperref}
\usepackage{setspace}
\usepackage{geometry}
\geometry{margin=1in}
\setstretch{1.3}

\title{\textbf{Testing and Evaluation of the Inventory Management System (IMS)}}
\author{Hruday Dessai}
\date{28/10/2025}

\begin{document}

\maketitle

\begin{abstract}
This paper presents the testing methodology and evaluation results of the \textbf{Inventory Management System (IMS)} project. The objective is to ensure that the system performs reliably under various operational conditions and meets user expectations for efficiency, scalability, and accuracy. Through comprehensive unit, integration, and performance testing, the IMS demonstrates a high level of robustness and stability suitable for real-world deployment.
\end{abstract}

\section{Introduction}
Testing plays a critical role in ensuring the overall quality, functionality, and reliability of any software system. The \textbf{Inventory Management System (IMS)} developed for automating inventory operations has undergone extensive testing across multiple levels.

The source code and documentation for the IMS are available on GitHub at:

\begin{center}
\href{https://github.com/NikhilChari/Inventory-Management-System.git}{https://github.com/NikhilChari/Inventory-Management-System.git}
\end{center}

This paper provides a detailed discussion on the testing methodology, evaluation metrics, and results, followed by an analysis of the challenges encountered and the improvements made.

\section{Testing Methodology}
A structured and multi-layered testing approach was followed to ensure the reliability and consistency of the IMS. The testing process was categorized into \textbf{Unit Testing}, \textbf{Integration Testing}, and \textbf{Performance Testing}.

\subsection*{1. Unit Testing}
Individual modules and functions were tested using frameworks such as \textbf{Jest} or \textbf{Mocha}. The goal was to verify that each component performed as expected in isolation.  
\textbf{Example modules tested:}
\begin{itemize}
    \item Inventory addition and update functions.
    \item User authentication and authorization modules.
    \item Database query and data validation routines.
\end{itemize}
The unit tests achieved an average code coverage of approximately \textbf{92\%}, ensuring that most functionalities were thoroughly validated.

\subsection*{2. Integration Testing}
Integration testing ensured seamless interaction between system components such as the frontend, backend APIs, and database layers.  
\textbf{Tools used:} \textbf{Postman} for API endpoint validation and \textbf{Supertest} for automated HTTP request testing.  
All RESTful APIs were validated for correct input/output behavior, authentication mechanisms, and error handling.

\subsection*{3. Performance and Load Testing}
To assess the scalability and responsiveness of the system, performance testing was carried out using tools like \textbf{Apache JMeter}.  
The system was tested under simulated loads with concurrent users performing database-intensive operations such as adding and querying stock records.  
Results indicated:
\begin{itemize}
    \item Average response time: \textbf{below 200 ms}.
    \item Peak concurrent users supported: \textbf{500}.
    \item Minimal latency during high-load scenarios.
\end{itemize}

\subsection*{4. User Acceptance Testing (UAT)}
A group of test users, including students and developers, evaluated the IMS interface and functionality. Feedback was collected regarding ease of use, visual clarity, and feature completeness.  
The overall \textbf{usability score} averaged \textbf{4.5 out of 5}, indicating high user satisfaction.

\section{Evaluation Metrics}
The following metrics were used to evaluate the IMS performance and effectiveness:

\begin{itemize}
    \item \textbf{Accuracy:} Correctness of operations, verified through cross-checking database updates with user actions.
    \item \textbf{Response Time:} Measured average latency during API calls and database queries.
    \item \textbf{Scalability:} Evaluated by gradually increasing the user load.
    \item \textbf{User Feedback:} Collected through surveys from 50 test users to assess usability and satisfaction.
\end{itemize}

The system consistently maintained data integrity and responsiveness, even under simulated peak conditions.

\section{Challenges and Improvements}
During the testing phase, certain challenges were identified and systematically resolved:
\begin{itemize}
    \item \textbf{Concurrency Issues:} Simultaneous database updates caused minor data synchronization delays. This was mitigated through optimized query handling and transaction management.
    \item \textbf{API Rate Limitations:} Load testing revealed bottlenecks under high API request rates. Solutions involved implementing caching and optimizing server middleware.
    \item \textbf{Browser Compatibility:} Minor UI inconsistencies across browsers were corrected through responsive design adjustments and cross-platform testing.
\end{itemize}

The improvements enhanced system stability and ensured consistent user experience across multiple platforms.

\section{Conclusion}
Comprehensive testing and evaluation confirmed that the \textbf{Inventory Management System (IMS)} performs efficiently, with high accuracy and stability. The combination of automated testing, user feedback, and performance validation provides strong assurance of the system’s reliability.

The testing process not only validated the functional correctness but also guided the optimization of system performance and scalability. Future work includes continuous integration (CI) testing pipelines and automated regression testing to support ongoing project updates.



\end{document}
